% GENERATED FILE. Edit canonical lessons, then run npm run generate:books.
% canonical-source-hash: fnv1a64:d253b6164928af3b
% canonical-lessons: AR-C181-ghassala, AR-C181-mukayyif, AR-C181-ulba, AR-C181-maghsala, AR-C181-raff

\chapter{Washing machine, Air conditioner, Tin, Sink, Shelf}
\label{ch:ar-washing-machine-air-conditioner-tin-sink-shelf}
\hlchaptermodality{181}

\begin{chapteropening}
\textbf{By the end of this chapter:} \emph{I can say a washing machine, an air conditioner, a tin, a sink and a shelf.}

Complete the last lesson of chapter 181: I can say a washing machine, an air conditioner, a tin, a sink and a shelf.
\end{chapteropening}

\section[\texorpdfstring{\ar{غسالة} (ghassāla)}{غسالة (ghassāla)}]{\ar{غسالة} (ghassāla) — a washing machine}

\label{lesson:AR-C181-ghassala}



\begin{quote}
Before the new one: say the Arabic for foreign, then the Arabic for flat.
\end{quote}

\subsection*{You'll want to know: \ar{غسالة}}
\textbf{\ar{غسالة}} — \emph{ghassāla} — \textquotedblleft{}a washing machine\textquotedblright{}.

\begin{cousinweb}
\textbf{\ar{غسالة}} (\emph{ghassāla}) is built on the native Arabic root \emph{gh-s-l}, \textquotedblleft{}to wash\textquotedblright{} (Lane). \emph{Ghassāla} is a modern \emph{faʿʿāla} word, \textquotedblleft{}the washer\textquotedblright{}.
\end{cousinweb}

\begin{grammarlens}[title={What you've built}]
One more word for this chapter.
\end{grammarlens}

\subsection*{Your turn}
\begin{itemize}
\raggedright
  \item \emph{Say it:} \emph{ghassāla}
  \item \emph{Say it:} \emph{ghassāla}, once more
  \item \emph{Say it:} say \emph{ghassāla}
  \item \emph{Recall:} say the Arabic for foreign, then the Arabic for flat, then say \emph{ghassāla} again
\end{itemize}

\begin{tcolorbox}[breakable,colback=teal!4,colframe=teal!35!black,title={Before you move on}]
What does \ar{غسالة} mean? (\textquotedblleft{}A washing machine\textquotedblright{}.) Say it once more.
\end{tcolorbox}
\section[\texorpdfstring{\ar{مكيف} (mukayyif)}{مكيف (mukayyif)}]{\ar{مكيف} (mukayyif) — an air conditioner}

\label{lesson:AR-C181-mukayyif}



\begin{quote}
Before the new one: say the Arabic for flat, then the Arabic for a washing machine.
\end{quote}

\subsection*{You'll want to know: \ar{مكيف}}
\textbf{\ar{مكيف}} — \emph{mukayyif} — \textquotedblleft{}an air conditioner\textquotedblright{}.

\begin{cousinweb}
\textbf{\ar{مكيف}} (\emph{mukayyif}) is built on \emph{k-y-f}, the letters of \emph{kayfa}, \textquotedblleft{}how\textquotedblright{} (Lane). \emph{Mukayyif} is a modern Form II participle, \textquotedblleft{}one that sets the state (of the air)\textquotedblright{}, built from \emph{kayfa}, \textquotedblleft{}how\textquotedblright{}.
\end{cousinweb}

\begin{grammarlens}[title={What you've built}]
One more word for this chapter.
\end{grammarlens}

\subsection*{Your turn}
\begin{itemize}
\raggedright
  \item \emph{Say it:} \emph{mukayyif}
  \item \emph{Say it:} \emph{mukayyif}, once more
  \item \emph{Say it:} say \emph{mukayyif}
  \item \emph{Recall:} say the Arabic for flat, then the Arabic for a washing machine, then say \emph{mukayyif} again
\end{itemize}

\begin{tcolorbox}[breakable,colback=teal!4,colframe=teal!35!black,title={Before you move on}]
What does \ar{مكيف} mean? (\textquotedblleft{}An air conditioner\textquotedblright{}.) Say it once more.
\end{tcolorbox}
\section[\texorpdfstring{\ar{علبة} (ʿulba)}{علبة (ʿulba)}]{\ar{علبة} (ʿulba) — a tin, a can}

\label{lesson:AR-C181-ulba}



\begin{quote}
Before the new one: say the Arabic for a washing machine, then the Arabic for an air conditioner.
\end{quote}

\subsection*{You'll want to know: \ar{علبة}}
\textbf{\ar{علبة}} — \emph{ʿulba} — \textquotedblleft{}a tin, a can\textquotedblright{}.

\begin{cousinweb}
\textbf{\ar{علبة}} (\emph{ʿulba}), \textquotedblleft{}a milking vessel\textquotedblright{}, is an old native noun, filed under the root \emph{ʿ-l-b} (Lane). Lane's \emph{ʿulba} is a milking vessel of skin or wood; the tin can is a modern sense.
\end{cousinweb}

\begin{grammarlens}[title={What you've built}]
One more word for this chapter.
\end{grammarlens}

\subsection*{Your turn}
\begin{itemize}
\raggedright
  \item \emph{Say it:} \emph{ʿulba}
  \item \emph{Say it:} \emph{ʿulba}, once more
  \item \emph{Say it:} say \emph{ʿulba}
  \item \emph{Recall:} say the Arabic for a washing machine, then the Arabic for an air conditioner, then say \emph{ʿulba} again
\end{itemize}

\begin{tcolorbox}[breakable,colback=teal!4,colframe=teal!35!black,title={Before you move on}]
What does \ar{علبة} mean? (\textquotedblleft{}A tin, a can\textquotedblright{}.) Say it once more.
\end{tcolorbox}
\section[\texorpdfstring{\ar{مغسلة} (maghsala)}{مغسلة (maghsala)}]{\ar{مغسلة} (maghsala) — a sink, a washbasin}

\label{lesson:AR-C181-maghsala}



\begin{quote}
Before the new one: say the Arabic for an air conditioner, then the Arabic for a tin.
\end{quote}

\subsection*{You'll want to know: \ar{مغسلة}}
\textbf{\ar{مغسلة}} — \emph{maghsala} — \textquotedblleft{}a sink, a washbasin\textquotedblright{}.

\begin{cousinweb}
\textbf{\ar{مغسلة}} (\emph{maghsala}) is built on the native Arabic root \emph{gh-s-l}, \textquotedblleft{}to wash\textquotedblright{} (Lane). \emph{Maghsala} is the \emph{ma-} place noun, \textquotedblleft{}a place of washing\textquotedblright{}.
\end{cousinweb}

\begin{grammarlens}[title={What you've built}]
One more word for this chapter.
\end{grammarlens}

\subsection*{Your turn}
\begin{itemize}
\raggedright
  \item \emph{Say it:} \emph{maghsala}
  \item \emph{Say it:} \emph{maghsala}, once more
  \item \emph{Say it:} say \emph{maghsala}
  \item \emph{Recall:} say the Arabic for an air conditioner, then the Arabic for a tin, then say \emph{maghsala} again
\end{itemize}

\begin{tcolorbox}[breakable,colback=teal!4,colframe=teal!35!black,title={Before you move on}]
What does \ar{مغسلة} mean? (\textquotedblleft{}A sink, a washbasin\textquotedblright{}.) Say it once more.
\end{tcolorbox}
\section[\texorpdfstring{\ar{رف} (raff)}{رف (raff)}]{\ar{رف} (raff) — a shelf}

\label{lesson:AR-C181-raff}



\begin{quote}
Before the new one: say the Arabic for a tin, then the Arabic for a sink.
\end{quote}

\subsection*{You'll want to know: \ar{رف}}
\textbf{\ar{رف}} — \emph{raff} — \textquotedblleft{}a shelf\textquotedblright{}.

\begin{cousinweb}
\textbf{\ar{رف}} (\emph{raff}), \textquotedblleft{}a shelf\textquotedblright{}, is an old native noun, filed under the root \emph{r-f-f} (Lane). Lane's \emph{raff} is a ledge on which the choice vessels of a house are set.
\end{cousinweb}

\begin{grammarlens}[title={What you've built}]
One more word for this chapter.
\end{grammarlens}

\subsection*{Your turn}
\begin{itemize}
\raggedright
  \item \emph{Say it:} \emph{raff}
  \item \emph{Say it:} \emph{raff}, once more
  \item \emph{Say it:} say \emph{raff}
  \item \emph{Recall:} say the Arabic for a tin, then the Arabic for a sink, then say \emph{raff} again
\end{itemize}

\begin{tcolorbox}[breakable,colback=teal!4,colframe=teal!35!black,title={Before you move on}]
What does \ar{رف} mean? (\textquotedblleft{}A shelf\textquotedblright{}.) Say it once more.
\end{tcolorbox}
